% ============================================================
\chapter{Carath\'eodory Extension Theorem}
\label{ch:caratheodory}
% ============================================================

\section{Introduction}

How does one \emph{construct} a measure? The question may seem circular: we want to measure sets, but to rigorously define what ``measuring'' means, we first need a theory. It was Constantin Carath\'eodory who, in 1914, cut this Gordian knot with a construction of remarkable elegance. His idea: start with a cruder notion, the \emph{outer measure}, which has the defect of not being additive in general, then select among all sets those on which this outer measure behaves well. The ``measurable'' sets in Carath\'eodory's sense automatically form a $\sigma$-algebra, and the outer measure restricted to this $\sigma$-algebra is a genuine measure.

This construction, as elegant as it is abstract, is the universal machine that produces the Lebesgue measure, the Hausdorff measure, and many others. It is the cornerstone on which the entire edifice of modern integration rests.

\section{Outer measures}

\begin{definition}[Outer measure]
Let $X$ be a set. An \emph{outer measure} on $X$ is a function
$\mu^* : \mathcal{P}(X) \to [0, +\infty]$ satisfying:
\begin{enumerate}
    \item[(i)] $\mu^*(\varnothing) = 0$;
    \item[(ii)] (monotonicity) if $A \subset B$, then $\mu^*(A) \leq \mu^*(B)$;
    \item[(iii)] (countable subadditivity) for every sequence
    $(A_n)_{n \in \N} \subset \mathcal{P}(X)$:
    \[
    \mu^*\!\left(\bigcup_{n=0}^{\infty} A_n\right) \leq \sum_{n=0}^{\infty}
    \mu^*(A_n).
    \]
\end{enumerate}
\end{definition}

\begin{remark}
An outer measure is defined on \emph{all} subsets of $X$, but it is generally not
$\sigma$-additive. Carath\'eodory's idea is to identify the sets on which it is
additive.
\end{remark}

\section{Carath\'eodory-measurable sets}

\begin{definition}[Carath\'eodory measurability]
\label{def:caratheodory_measurable}
Let $\mu^*$ be an outer measure on $X$. A set $A \subset X$ is
\emph{$\mu^*$-measurable} (or \emph{Carath\'eodory-measurable}) if:
\[
\forall\, E \subset X, \quad \mu^*(E) = \mu^*(E \cap A) + \mu^*(E \cap A^c).
\]
We denote by $\mathcal{M}^*$ the collection of all $\mu^*$-measurable sets.
\end{definition}

\begin{intuition}[title=\textbf{Interpreting the Carath\'eodory criterion}]
The condition says that $A$ ``splits'' every set $E$ properly: the outer measure
of $E$ equals the sum of the outer measures of the two pieces $E \cap A$ and
$E \cap A^c$. By subadditivity, the inequality
$\mu^*(E) \leq \mu^*(E \cap A) + \mu^*(E \cap A^c)$ always holds. So it suffices
to verify the reverse inequality:
$\mu^*(E) \geq \mu^*(E \cap A) + \mu^*(E \cap A^c)$.
\end{intuition}

\section{The main theorem}

\begin{theorem}[Carath\'eodory]
\label{thm:caratheodory}
Let $\mu^*$ be an outer measure on $X$. Then:
\begin{enumerate}
    \item $\mathcal{M}^*$ is a $\sigma$-algebra;
    \item the restriction $\mu = \mu^*|_{\mathcal{M}^*}$ is a complete measure on
    $(X, \mathcal{M}^*)$.
\end{enumerate}
\end{theorem}

\begin{proof}
The proof proceeds in several steps.

\emph{Step 1: $\mathcal{M}^*$ is an algebra.}

Clearly $X \in \mathcal{M}^*$ (take $A = X$: $\mu^*(E) = \mu^*(E) +
\mu^*(\varnothing) = \mu^*(E)$). If $A \in \mathcal{M}^*$, then $A^c \in
\mathcal{M}^*$ since the Carath\'eodory condition is symmetric in $A$ and $A^c$.

Let us show closure under finite unions. Let $A, B \in \mathcal{M}^*$. For every
$E \subset X$:
\begin{align*}
\mu^*(E) &= \mu^*(E \cap A) + \mu^*(E \cap A^c) \\
&= \mu^*(E \cap A \cap B) + \mu^*(E \cap A \cap B^c) + \mu^*(E \cap A^c \cap B)
+ \mu^*(E \cap A^c \cap B^c).
\end{align*}
Now $E \cap (A \cup B) = (E \cap A \cap B) \cup (E \cap A \cap B^c) \cup
(E \cap A^c \cap B)$ and $E \cap (A \cup B)^c = E \cap A^c \cap B^c$. By
subadditivity:
\[
\mu^*(E \cap (A \cup B)) \leq \mu^*(E \cap A \cap B) + \mu^*(E \cap A \cap B^c)
+ \mu^*(E \cap A^c \cap B).
\]
Therefore $\mu^*(E) \geq \mu^*(E \cap (A \cup B)) + \mu^*(E \cap (A \cup B)^c)$,
and the reverse inequality follows from subadditivity. Thus $A \cup B \in
\mathcal{M}^*$.

\emph{Step 2: Finite additivity.}

If $A, B \in \mathcal{M}^*$ are disjoint, taking $E = A \cup B$:
$\mu^*(A \cup B) = \mu^*((A \cup B) \cap A) + \mu^*((A \cup B) \cap A^c)
= \mu^*(A) + \mu^*(B)$.

By induction, for $A_1, \ldots, A_n \in \mathcal{M}^*$ pairwise disjoint:
$\mu^*\!\left(\bigsqcup_{k=1}^n A_k\right) = \sum_{k=1}^n \mu^*(A_k)$.

\emph{Step 3: $\sigma$-additivity and closure under countable unions.}

Let $(A_n)_{n \geq 1} \subset \mathcal{M}^*$ be a disjoint sequence. Set
$S_n = \bigsqcup_{k=1}^n A_k$ and $A = \bigsqcup_{k=1}^\infty A_k$. By Step 1,
$S_n \in \mathcal{M}^*$, so for every $E$:
\[
\mu^*(E) = \mu^*(E \cap S_n) + \mu^*(E \cap S_n^c)
\geq \sum_{k=1}^n \mu^*(E \cap A_k) + \mu^*(E \cap A^c)
\]
since $S_n^c \supset A^c$. Letting $n \to \infty$:
\[
\mu^*(E) \geq \sum_{k=1}^{\infty} \mu^*(E \cap A_k) + \mu^*(E \cap A^c)
\geq \mu^*(E \cap A) + \mu^*(E \cap A^c)
\]
by subadditivity of $\mu^*$. The reverse inequality always holds, so
$A \in \mathcal{M}^*$.

Taking $E = A$: $\mu^*(A) = \sum_{k=1}^{\infty} \mu^*(A_k) + \mu^*(\varnothing)
= \sum_k \mu^*(A_k)$, giving $\sigma$-additivity.

\emph{Step 4: Completeness.}

Let $N \in \mathcal{M}^*$ with $\mu^*(N) = 0$ and $A \subset N$. For every
$E \subset X$: $\mu^*(E \cap A) \leq \mu^*(N) = 0$ and $\mu^*(E \cap A^c) \leq
\mu^*(E)$, so $\mu^*(E \cap A) + \mu^*(E \cap A^c) \leq \mu^*(E)$. The reverse
inequality is subadditivity, so $A \in \mathcal{M}^*$.
\end{proof}

\section{Application to pre-measures}

\begin{theorem}[Carath\'eodory extension]
\label{thm:caratheodory_extension}
Let $\mu_0$ be a pre-measure on an algebra $\mathcal{A}_0 \subset \mathcal{P}(X)$.
Define:
\[
\mu^*(A) = \inf\left\{\sum_{n=0}^{\infty} \mu_0(A_n) :
(A_n) \subset \mathcal{A}_0,\, A \subset \bigcup_{n=0}^{\infty} A_n\right\},
\quad A \subset X.
\]
Then:
\begin{enumerate}
    \item $\mu^*$ is an outer measure on $X$.
    \item $\mathcal{A}_0 \subset \mathcal{M}^*$ (every element of the algebra is
    Carath\'eodory-measurable).
    \item $\mu^*|_{\mathcal{A}_0} = \mu_0$ (the extension agrees with the
    pre-measure on the algebra).
    \item In particular, $\mu = \mu^*|_{\sigma(\mathcal{A}_0)}$ is a measure on
    $\sigma(\mathcal{A}_0)$ extending $\mu_0$.
\end{enumerate}
\end{theorem}

\begin{proof}
\emph{Part 1.} $\mu^*(\varnothing) = 0$ (take $A_n = \varnothing$). Monotonicity
is clear. For countable subadditivity, let $(B_k)$ be a sequence and
$\varepsilon > 0$. For each $k$, choose $(A_n^{(k)}) \subset \mathcal{A}_0$ with
$B_k \subset \bigcup_n A_n^{(k)}$ and $\sum_n \mu_0(A_n^{(k)}) \leq
\mu^*(B_k) + \varepsilon/2^{k+1}$. Then $\bigcup_k B_k \subset
\bigcup_{k,n} A_n^{(k)}$ and
\[
\mu^*\!\left(\bigcup_k B_k\right) \leq \sum_{k,n} \mu_0(A_n^{(k)}) \leq
\sum_k \mu^*(B_k) + \varepsilon.
\]
Letting $\varepsilon \to 0$ gives countable subadditivity.

\emph{Part 2.} Let $A \in \mathcal{A}_0$ and $E \subset X$. It suffices to show
$\mu^*(E) \geq \mu^*(E \cap A) + \mu^*(E \cap A^c)$. Let $\varepsilon > 0$ and
$(A_n) \subset \mathcal{A}_0$ with $E \subset \bigcup_n A_n$ and
$\sum_n \mu_0(A_n) \leq \mu^*(E) + \varepsilon$. Since $\mathcal{A}_0$ is an
algebra, $A_n \cap A$ and $A_n \cap A^c$ are in $\mathcal{A}_0$ and
$\mu_0(A_n) = \mu_0(A_n \cap A) + \mu_0(A_n \cap A^c)$ (finite additivity). Thus:
\begin{align*}
\mu^*(E) + \varepsilon &\geq \sum_n \mu_0(A_n \cap A) + \sum_n \mu_0(A_n \cap A^c)
\\ &\geq \mu^*(E \cap A) + \mu^*(E \cap A^c).
\end{align*}

\emph{Part 3.} For $A \in \mathcal{A}_0$, $\mu^*(A) \leq \mu_0(A)$ (take
$A_0 = A$, $A_n = \varnothing$). For the reverse, let $(A_n) \subset \mathcal{A}_0$
with $A \subset \bigcup_n A_n$. Set $B_N = A \cap \bigl(\bigcup_{n=0}^N A_n\bigr)
\in \mathcal{A}_0$. Then $B_N \uparrow A$ and $B_N \subset \bigcup_{n=0}^N A_n$.
By finite subadditivity: $\mu_0(B_N) \leq \sum_{n=0}^N \mu_0(A_n)$. By
continuity from below (Proposition~\ref{prop:sigma_add_criterion}):
$\mu_0(A) = \lim_N \mu_0(B_N) \leq \sum_n \mu_0(A_n)$. Taking the infimum:
$\mu_0(A) \leq \mu^*(A)$.
\end{proof}

\section{Uniqueness of the extension}

\begin{theorem}[Uniqueness]
\label{thm:uniqueness_extension}
If the pre-measure $\mu_0$ is $\sigma$-finite on $\mathcal{A}_0$, then the extension
$\mu$ to $\sigma(\mathcal{A}_0)$ is unique.
\end{theorem}

\begin{proof}
By the uniqueness theorem for measures (Theorem~\ref{thm:uniqueness_measures}),
applied to the $\pi$-system $\mathcal{A}_0$ (an algebra is a $\pi$-system). The
$\sigma$-finiteness of $\mu_0$ provides the required exhausting sequence.
\end{proof}

\begin{keyformulas}[title=\textbf{Summary of the Carath\'eodory construction}]
\[
\text{Pre-measure } \mu_0 \text{ on } \mathcal{A}_0
\xrightarrow{\text{outer meas.}} \mu^* \text{ on } \mathcal{P}(X)
\xrightarrow{\text{restrict}} \mu \text{ on } \mathcal{M}^*
\supset \sigma(\mathcal{A}_0).
\]
\end{keyformulas}

\section{Carath\'eodory criterion for metric outer measures}

\begin{definition}[Metric outer measure]
Let $(X, d)$ be a metric space. An outer measure $\mu^*$ on $X$ is \emph{metric} if:
\[
d(A, B) > 0 \implies \mu^*(A \cup B) = \mu^*(A) + \mu^*(B),
\]
where $d(A, B) = \inf\{d(a, b) : a \in A,\, b \in B\}$.
\end{definition}

\begin{theorem}
If $\mu^*$ is a metric outer measure on $(X, d)$, then every Borel set is
$\mu^*$-measurable: $\mathcal{B}(X) \subset \mathcal{M}^*$.
\end{theorem}

\begin{proof}
It suffices to show that every closed set $F$ is $\mu^*$-measurable. Let
$E \subset X$ and set $E_n = \{x \in E \cap F^c : d(x, F) \geq 1/n\}$. Then
$E_n \uparrow E \cap F^c$ and $d(E \cap F, E_n) \geq 1/n > 0$, so
$\mu^*(E \cap F \cup E_n) = \mu^*(E \cap F) + \mu^*(E_n)$ by the metric property.

By monotonicity, $\mu^*(E) \geq \mu^*(E \cap F \cup E_n) = \mu^*(E \cap F) +
\mu^*(E_n)$. It remains to show $\mu^*(E_n) \to \mu^*(E \cap F^c)$.

Set $D_n = E_{n+1} \setminus E_n$. For $n$ and $m$ with $m \geq n + 2$,
$d(D_n, D_m) > 0$, so the sets $D_n$ of fixed parity are at positive distance from
each other. By the metric property and countable subadditivity:
\[
\mu^*(E \cap F^c) \leq \mu^*(E_n) + \sum_{k=n}^{\infty} \mu^*(D_k).
\]
If $\sum_k \mu^*(D_k) < \infty$, the tail tends to zero and
$\mu^*(E_n) \to \mu^*(E \cap F^c)$. If the series diverges, one shows that
$\mu^*(E) = +\infty$ and the inequality is trivial.
\end{proof}

\section{Application: Hausdorff measure}

\begin{definition}[Hausdorff measure]
Let $(X, d)$ be a metric space and $s \geq 0$. For $\delta > 0$, define:
\[
\mathcal{H}^s_\delta(A) = \inf\left\{\sum_{n=1}^{\infty} (\operatorname{diam}
U_n)^s : A \subset \bigcup_{n=1}^{\infty} U_n,\,
\operatorname{diam} U_n \leq \delta\right\}.
\]
The \emph{$s$-dimensional Hausdorff measure} is:
\[
\mathcal{H}^s(A) = \lim_{\delta \to 0} \mathcal{H}^s_\delta(A) =
\sup_{\delta > 0} \mathcal{H}^s_\delta(A).
\]
\end{definition}

\begin{proposition}
$\mathcal{H}^s$ is a metric outer measure. Consequently, it is a (Borel) measure
on $(X, \mathcal{B}(X))$.
\end{proposition}

\section{Exercises}

\begin{exercise}
Verify in detail that the set $\mathcal{M}^*$ in Carath\'eodory's theorem is closed
under countable intersections.
\end{exercise}

\begin{exercise}
Let $\mu^*$ be an outer measure on $X$. Show that if $\mu^*(A) = 0$, then
$A \in \mathcal{M}^*$.
\end{exercise}

\begin{exercise}
Show that the Hausdorff measure $\mathcal{H}^n$ on $\R^n$ equals (up to a
multiplicative constant) the Lebesgue measure.
\end{exercise}

\begin{exercise}
Let $\mu_0$ be a \emph{finite} pre-measure on an algebra $\mathcal{A}_0$. Show that
for every $A \in \sigma(\mathcal{A}_0)$ and every $\varepsilon > 0$, there exists
$B \in \mathcal{A}_0$ such that $\mu(A \triangle B) < \varepsilon$, where
$A \triangle B = (A \setminus B) \cup (B \setminus A)$.
\end{exercise}

\begin{exercise}
Let $X$ be an uncountable set and
$\mu^*(A) = \begin{cases} 0 & \text{if } A \text{ is countable,} \\
1 & \text{otherwise.} \end{cases}$
Show that $\mu^*$ is an outer measure and determine $\mathcal{M}^*$.
\end{exercise}
